\documentclass[11pt,aspectratio=169]{beamer}
\usetheme{Boadilla}
\usefonttheme{professionalfonts}
\setbeamertemplate{navigation symbols}{}

\usepackage[T1]{fontenc}
\usepackage{lmodern}
\usepackage{amsmath,amssymb,mathtools}
\usepackage{booktabs}
\usepackage{tabularx}
\usepackage{array}
\usepackage{graphicx}
\usepackage[percent]{overpic}
\usepackage{xcolor}
\usepackage[authoryear,round]{natbib}
\usepackage{bibentry}
\bibliographystyle{apalike}
\setbeamertemplate{bibliography item}{}
\makeatletter
\def\hyper@natlinkstart#1{}
\def\hyper@natlinkend{}
\def\hyper@natlinkbreak#1#2{#1}
\makeatother

\graphicspath{{figs/}}

% Section headline for the audience arc.
\setbeamertemplate{headline}{%
  \hbox{}\vskip2pt\hfill\textcolor{gray!70}{\scriptsize\insertsectionhead}\hspace{1em}\vskip-6pt%
}
\setbeamertemplate{section in toc}[sections numbered]
\setbeamertemplate{frametitle}[default][left]
\setbeamerfont{framesubtitle}{size=\small}

\newcommand{\figplaceholder}[1]{\begin{center}\fbox{\begin{minipage}[c][0.32\textheight][c]{0.82\textwidth}\centering\large #1\end{minipage}}\end{center}}
\newcommand{\shortfigplaceholder}[1]{\begin{center}\fbox{\begin{minipage}[c][0.22\textheight][c]{0.82\textwidth}\centering\large #1\end{minipage}}\end{center}}
\newcommand{\compactfigplaceholder}[1]{\begin{center}\fbox{\begin{minipage}[c][0.10\textheight][c]{0.82\textwidth}\centering\normalsize #1\end{minipage}}\end{center}}
\newcommand{\R}{\mathcal{R}}
\newcommand{\E}{\mathcal{E}}
\newcommand{\dd}{\,\mathrm{d}}
\newcommand{\good}[1]{\textcolor{blue!70!black}{#1}}
\newcommand{\bad}[1]{\textcolor{red!70!black}{#1}}
\newcommand{\SAR}{speed asymmetry ratio}
% Inline slide navigation (end of sentence / bullet).
\newcommand{\app}[1]{\allowbreak\hspace{0.25em}\hyperlink{#1}{\beamergotobutton{details}}}
\newcommand{\back}[1]{\allowbreak\hspace{0.25em}\hyperlink{#1}{\beamerreturnbutton{back}}}
% Appendix / frame footer: bottom-right, not centered.
\newcommand{\navback}[1]{\par\vspace{0.35em}\noindent\hfill\hyperlink{#1}{\beamerreturnbutton{back}}}
\newcommand{\navfwd}[1]{\par\vspace{0.35em}\noindent\hfill\hyperlink{#1}{\beamergotobutton{details}}}
\newcommand{\navbackfwd}[2]{\par\vspace{0.35em}\noindent\hfill\hyperlink{#1}{\beamerreturnbutton{back}}\;\hyperlink{#2}{\beamergotobutton{details}}}

\title[Directional Congestion]{Directional Congestion}
\subtitle{Measurement, Spatial Equilibrium, and Tidal-Lane Reallocation}
\author[Duan, Feng, Gu]{Yige Duan (SJTU) $\quad$ Xiaoruo Feng (PKU) $\quad$ Yizhen Gu (PKU)}
\date[SMUE Wuhan, 2026]{SMUE Wuhan, 2026}

\begin{document}

\begin{frame}
\titlepage
\end{frame}

% =====================================================================
% Opener: roadmap followed by examples of directional congestion.
% =====================================================================
\begin{frame}{Roadmap}
\footnotesize\tableofcontents
\end{frame}

\begin{frame}{Directional Congestion in Cities}

\centering
% All four photos center-cropped to 3:2 for a uniform grid.
\begin{minipage}[t]{0.32\textwidth}\centering
\includegraphics[width=0.8\linewidth]{cityBeijing_c.jpg}\\[0.02em]{\scriptsize Beijing}
\end{minipage}\hspace{0.02\textwidth}%
\begin{minipage}[t]{0.32\textwidth}\centering
\includegraphics[width=0.8\linewidth]{cityWuhan_c.jpg}\\[0.02em]{\scriptsize Wuhan}
\end{minipage}

\vspace{0.25em}
\begin{minipage}[t]{0.32\textwidth}\centering
\includegraphics[width=0.8\linewidth]{cityShenzhen2_c.jpg}\\[0.02em]{\scriptsize Shenzhen}
\end{minipage}\hspace{0.02\textwidth}%
\begin{minipage}[t]{0.32\textwidth}\centering
\includegraphics[width=0.8\linewidth]{cityKualaLumpur_c.jpg}\\[0.02em]{\scriptsize Kuala Lumpur}
\end{minipage}
\end{frame}

% =====================================================================
\section{Introduction}
% =====================================================================
\begin{frame}{Motivation}
\begin{columns}[T,onlytextwidth]
\column{0.65\textwidth}
\begin{itemize}\small
  \item \textbf{Directional congestion is widely observed on urban roads, particularly during peak hours.}
  \begin{itemize}
    \item Commuting traffic is shaped by where workers live and where jobs are located.
  \end{itemize}
  \item \textbf{The literature often ignores this asymmetry in congestion.}
  \begin{itemize}
    \item Examples: QSM applications usually enter travel frictions through bilateral or route-level costs \citep{ahlfeldt2015,allenarkolakis2022}.
    \item Except: transport studies discuss peak-direction traffic \citep{vickrey1969,wolshon2006}.
  \end{itemize}
  \item \textbf{Tidal lanes are proposed as a low-cost way to address it.}
  \begin{itemize}
    \item However, little is known about the short-run and long-run effects of tidal lanes.
  \end{itemize}
\end{itemize}
\column{0.34\textwidth}
\centering
\href{run:./A6_XiaoruoFeng_DirectionalCongestion_video.mov}{%
  \includegraphics[width=\linewidth,height=0.70\textheight,keepaspectratio]{TidalVedio/pic.jpg}}\\[3pt]
{\scriptsize \good{Click to play}}
\end{columns}
\end{frame}

\begin{frame}{This Paper}
\small
\begin{itemize}
  \item \textbf{Measure directionality, using speed data at both the trip and the road segment level.}
  \begin{itemize}\footnotesize
    \item National trip-level: flow-weighted median \SAR\ is \good{0.878}.
    \item Within-city road segments: \SAR\ $\leq1/3$ flags \good{179} Beijing AM two-way centerlines (\good{4.9\%}, about \good{213 km}).
  \end{itemize}
  \item \textbf{Quantify what symmetry misses, using a spatial GE model.}
  \begin{itemize}\footnotesize
    \item Holding fundamentals and parameters fixed, replacing directed costs with reciprocal symmetric costs changes $\bar W$ by \good{$1.10\%$} and moves \good{$3.36\%$} of residents and \good{$6.15\%$} of jobs.
  \end{itemize}
  \item \textbf{Evaluate counterfactuals under endogenous congestion GE.}
  \begin{itemize}\footnotesize
    \item Tidal top-100 shifts one lane toward the slow AM direction on each treated corridor.
    \item Citywide physical time falls by \good{$42.1$ seconds per commuter} in ST; LT reallocation raises it by \bad{$14.4$ seconds} above baseline.
    \item Welfare rises \good{$+0.87\%$} in ST and \good{$+0.32\%$} in LT.
    \item Capacity-expansion benchmark builds an additional \good{$1{,}100$ lane-km}, while tidal lanes add \good{zero}.
  \end{itemize}
\end{itemize}
\end{frame}

\begin{frame}{Literature}
\begin{itemize}\small
  \item \textbf{Directional congestion}
  \begin{itemize}\footnotesize
    \item Peak-period bottlenecks, measured urban speeds, and reversible-lane operations \citep{vickrey1969,arnott1993,couture2018,akbar2023,wolshon2006}.
    \item \good{This paper measures directionality at national trip level and road-segment level using high-frequency speed data.}
  \end{itemize}
  \item \textbf{Spatial equilibrium}
  \begin{itemize}\footnotesize
    \item Economic geography and quantitative spatial models \citep{krugman1991,eaton2002,allen2014,ahlfeldt2015}.
    \item \good{This paper applies spatial GE with directional congestion and analyze tidal lane counterfactuals}
  \end{itemize}
  \item \textbf{Transportation policy and congestion}
  \begin{itemize}\footnotesize
    \item Demand-side policies: congestion pricing and equilibrium sorting \citep{vickrey1969,kreindler2024,barwick2024,yang2020}.
    \item Supply-side policies: road capacity, fundamental law, and optimal networks \citep{duranton2011,fajgelbaum2020}.
    \item Public transit as a substitute: subways relieve road congestion \citep{gu2021}.
    \item \good{This paper evaluates tidal lanes as a directed capacity-reallocation policy.}
  \end{itemize}
\end{itemize}
\end{frame}

\begin{frame}[label=nattrip]{Directionality Appears at National and Road-Segment Scales}
\begin{itemize}\footnotesize
  \item \textbf{Definition.} For each pair, the \SAR\ is the lower speed divided by the higher speed across its two directions. It lies between 0 and 1; lower values indicate stronger directional asymmetry.
\end{itemize}
\vspace{-0.35em}
\begin{columns}[T,onlytextwidth]
\column{0.60\textwidth}
\centering\scriptsize \textbf{National trip-level directionality}\par
\vspace{0.15em}
\begin{minipage}[t]{0.49\linewidth}\centering
\includegraphics[width=\linewidth,height=0.31\textheight,keepaspectratio]{pat_national_ratio_hist.png}
\end{minipage}\hspace{0.01\linewidth}%
\begin{minipage}[t]{0.49\linewidth}\centering
\includegraphics[width=\linewidth,height=0.31\textheight,keepaspectratio]{pat_congestion_am_scatter.png}
\end{minipage}
\vspace{-0.35em}
\begin{itemize}\scriptsize
  \item Flow-weighted median \SAR\ is \good{0.878} across the top 100 congested cities.
  \item \good{58\%} of pair-flow faces \SAR\ $\leq0.90$ during the AM peak.
  \item More congested cities tend to have lower ratios and stronger asymmetry ($r=\bad{-0.48}$, $n=99$).
\end{itemize}
\column{0.38\textwidth}
\centering\scriptsize \textbf{Within-city road-segment directionality}\par
\vspace{0.15em}
\includegraphics[width=\linewidth,height=0.42\textheight,keepaspectratio]{pat_bj_centerline_map.png}
\vspace{-0.15em}
\begin{itemize}\scriptsize
  \item Median \SAR\ is \good{0.884} among 3{,}670 Beijing two-way centerlines.
  \item \SAR\ $\leq1/3$ flags \good{179} candidate corridors, about \good{213 km}.
\end{itemize}
\end{columns}
\end{frame}

\begin{frame}[label=jobhousing]{Jobs--Housing Separation Generates Directional Commuting}
\begin{columns}[T,onlytextwidth]
\column{0.48\textwidth}
\centering
\small\textbf{Across major cities}\par
\vspace{0.30em}
\begin{itemize}\footnotesize\setlength\itemsep{0.55em}
  \item Core towns $C$ cover the first 30\% of workplace-side commuter mass; $P$ denotes the remaining towns.
  \item In all top-20 cities, $P\to C$ exceeds $C\to P$; the flow-weighted gap is \good{4.5 pp}.
\end{itemize}
\vspace{0.35em}
\includegraphics[width=0.96\linewidth,height=0.34\textheight,keepaspectratio]{pat_jobs_housing_top20_gap_column.pdf}

\column{0.50\textwidth}
\centering
\small\textbf{Beijing}\par
\vspace{0.30em}
\begin{itemize}\footnotesize\setlength\itemsep{0.55em}
  \item \textbf{Separation.} Jobs are more concentrated than residences: Duncan $D=\good{0.26}$ versus $0.004$ under random allocation. \app{appduncan}
  \item \textbf{Direction.} \good{74.5\%} of AM trips are center-bound; inbound/reverse is $\good{2.9\times}$. \app{appinbound}
\end{itemize}
\vspace{0.35em}
\includegraphics[width=0.96\linewidth,height=0.34\textheight,keepaspectratio]{pat_jobs_housing_balance.png}
\end{columns}
\end{frame}

% =====================================================================
\section{Model}
% =====================================================================
% Locked notation: m_kl is physical link time in minutes, t_kl is the
% link-level iceberg cost, and tau_ij is the all-route bilateral cost index.

\begin{frame}{Model: Environment}
\begin{itemize}\small
  \item Locations $i,j\in N$ host a fixed mass $\bar L$ of residents and jobs.
  \item A physical road is represented by two directed links, $(k,l)$ and $(l,k)$.
  \begin{itemize}\footnotesize
    \item $m_{kl}^{ff}$ is free-flow physical time, measured in minutes.
    \item $n_{kl}$ is the number of lanes assigned to that direction.
  \end{itemize}
  \item Realized physical time $m_{kl}$ is endogenous and directed.
  \begin{itemize}\footnotesize
    \item $m_{kl}$ rises with traffic per lane, $\Xi_{kl}/n_{kl}$.
    \item $m_{kl}\neq m_{lk}$ when traffic differs across directions.
  \end{itemize}
  \item The model maps minutes into link iceberg cost: $t_{kl}=m_{kl}^{1/\theta}$.
  \item A tidal-lane policy shifts one lane across reciprocal directions while holding $n_{kl}+n_{lk}$ fixed.
\end{itemize}
\end{frame}
\note{\textbf{Talk track, 35 seconds.} A physical road becomes two directed links. The data deliver physical minutes $m_{kl}$ and directional lanes. The model converts minutes into the link iceberg cost $t_{kl}=m_{kl}^{1/\theta}$. A tidal lane changes directed capacity without expanding total road width.}

\begin{frame}{Individual Location Choice}
\begin{itemize}\small
  \item An individual $\omega$ living in $i$, working in $j$, and taking route $r\in\R_{ij}$ obtains
  \vspace{-0.5em}\begin{footnotesize}\begin{equation}\label{eq:routecost}
    V_{ij,r}(\omega) = \left( \frac{u_i w_j}{C_{ij,r}} \right)\varepsilon_{ij,r}(\omega),
    \qquad
    C_{ij,r} = \prod_{\ell=1}^{K_r} t_{r_{\ell-1},r_\ell}
  \vspace{-1.25em}\end{equation}\end{footnotesize}
  \begin{itemize}\footnotesize
    \item $u_i$ is amenity; $w_j$ is wage; $\varepsilon_{ij,r}$ is i.i.d. Fr\'echet.
    \item $C_{ij,r}$ is route-specific traversal cost for $r=[r_0,\ldots,r_{K_r}]$, with $r_0=i$ and $r_{K_r}=j$.
  \end{itemize}
  \vspace{0.25em}
  \item Residence, workplace, and route are chosen jointly:
  \vspace{-0.5em}\begin{equation}\label{eq:pi}
    \pi_{ij,r} = \frac{\left(u_i w_j/C_{ij,r}\right)^\theta}{\sum_{a}\sum_b\sum_{s\in\R_{ab}}\left(u_a w_b/C_{ab,s}\right)^\theta}.
  \end{equation}
\end{itemize}
\end{frame}
\note{\textbf{Talk track, 45 seconds.} The choice object is $(i,j,r)$. Physical minutes add along a realized route, but iceberg costs compound multiplicatively. Fr\'echet heterogeneity delivers a tractable probability for every residence, workplace, and route combination.}

\begin{frame}[label=bilateral]{Bilateral Commuting Costs and Flows}
\begin{itemize}\small
  \item Aggregating route costs gives the bilateral commuting cost:
  \vspace{-0.5em}\begin{equation}\label{eq:bilateral}
    \tau_{ij} = \left[\sum_{r\in\R_{ij}} C_{ij,r}^{-\theta}\right]^{-1/\theta}.
  \end{equation}
  \vspace{-0.5em}
  \item Summing route-specific flows gives residence-workplace commuting gravity:
  \vspace{-0.35em}\begin{equation}\label{eq:gravity}
    L_{ij}=\sum_{r\in\R_{ij}}L_{ij,r}
    =\tau_{ij}^{-\theta}\left(\frac{L_i^R}{\Pi_i}\right)\left(\frac{L_j^F}{P_j}\right),
    \quad L_i^R=\sum_jL_{ij},\quad L_j^F=\sum_iL_{ij}.
  \vspace{-0.8em}\end{equation}
  \begin{itemize}\footnotesize
    \item $\Pi_i$ and $P_j$ are origin and destination market-access terms.
    \item Directed link costs imply $\tau_{ij}\neq\tau_{ji}$ in general.
    \item $\tau_{ij}$ is an all-route composite cost index, not OD travel minutes.
  \end{itemize}
\end{itemize}
\end{frame}
\note{\textbf{Talk track, 45 seconds.} The model first assigns probability to each route. Summing over routes produces one bilateral cost $\tau_{ij}$ and total OD flow $L_{ij}$. This gravity equation is route-inclusive because every directed link enters through $C_{ij,r}$.}

\begin{frame}{Amenity, Productivity, and Location Choices}
\begin{itemize}\small
  \item Local outcomes combine fundamentals with endogenous density:
  \vspace{-0.5em}\begin{equation}\label{eq:agglom}
    w_j = A_j = \bar A_j\left(L_j^F\right)^\alpha,
    \qquad
    u_i=\bar u_i\left(L_i^R\right)^\beta.
  \end{equation}
  \vspace{-0.4em}
  \item Residence ($R$) and workplace ($F$) shares satisfy
  \vspace{-0.2em}
  \begin{align}
  (l_i^R)^{1-\beta\theta} & = \chi \sum_j\tau_{ij}^{-\theta}\,\bar u_i^{\theta}\,\bar A_j^{\theta}\,(l_j^F)^{\alpha\theta} \label{eq:clear1} \\
  (l_j^F)^{1-\alpha\theta} & = \chi \sum_i\tau_{ij}^{-\theta}\,\bar u_i^{\theta}\,\bar A_j^{\theta}\,(l_i^R)^{\beta\theta}. \label{eq:clear2}
  \end{align}
  \vspace{-0.7em}
  \item $\alpha>0$ governs workplace agglomeration; $\beta<0$ governs residential crowding; $\chi$ normalizes shares to one.
\end{itemize}
\end{frame}
\note{\textbf{Talk track, 45 seconds.} Fundamentals are $\bar u_i$ and $\bar A_j$. Employment density affects productivity through $\alpha$, and residential density affects amenities through $\beta$. Bilateral commuting costs connect these two location distributions.}

\begin{frame}{Feedback between Location Choice, Route Traffic, and Traversal Cost}
\begin{itemize}\footnotesize
  \item For a directed trip $i\to j$, $\kappa^{kl}_{ij}$ is expected use of link $(k,l)$:
  \vspace{-0.8em}\begin{equation}
    \kappa^{kl}_{ij}=\left(\frac{\tau_{ij}}{\tau_{ik}\,t_{kl}\,\tau_{lj}}\right)^\theta.
  \vspace{-0.6em}\end{equation}
  \item Directed link traffic aggregates OD flows across routes:
  \vspace{-0.5em}\begin{equation}\label{eq:Xi_c}
    \Xi_{kl}=\sum_i\sum_j\kappa^{kl}_{ij}L_{ij}.
  \vspace{-0.6em}\end{equation}
  \item Traffic per lane raises physical link time, which maps into link cost:
  \vspace{-0.5em}\begin{equation}\label{eq:c_Xi}
    m_{kl}=m_{kl}^{ff}\left(\frac{\Xi_{kl}}{n_{kl}}\right)^\lambda,
    \qquad t_{kl}=m_{kl}^{1/\theta}.
  \vspace{-0.6em}\end{equation}
  \item Expected welfare is
  \vspace{-0.5em}\begin{equation}
    \bar W=\left[\sum_{i,j}\tau_{ij}^{-\theta}u_i^\theta w_j^\theta\right]^{1/\theta}.
  \vspace{-0.8em}\end{equation}
\end{itemize}
\end{frame}
\note{\textbf{Talk track, 60 seconds.} $\kappa_{ij}^{kl}$ maps one OD commuter into expected passages over a directed link. Summing over OD flows gives model-implied link traffic $\Xi_{kl}$. Traffic raises physical minutes $m_{kl}$, which are converted into $t_{kl}$ and feed back into location and route choices. This fixed point is the endogenous-congestion channel.}

\begin{frame}[label=equilibrium]{Equilibrium}
\begin{itemize}\footnotesize
  \item Given fundamentals $\{\bar u_i,\bar A_i,m_{kl}^{ff},n_{kl}\}$ and parameters $(\theta,\alpha,\beta,\lambda)$, an equilibrium is
  \vspace{-0.45em}\begin{equation}
    \left\{L_i^R,L_j^F,L_{ij},m_{kl},t_{kl},\tau_{ij},\Xi_{kl},\Pi_i,P_j\right\}
  \vspace{-0.55em}\end{equation}
  such that
  \begin{itemize}\scriptsize
    \item residence-workplace-route choices satisfy \eqref{eq:pi};
    \item bilateral flows satisfy gravity \eqref{eq:gravity} and location shares satisfy \eqref{eq:clear1}--\eqref{eq:clear2};
    \item amenities and productivity satisfy \eqref{eq:agglom};
    \item bilateral costs aggregate route-specific directed link costs through \eqref{eq:routecost} and \eqref{eq:bilateral};
    \item directed link traffic and traversal costs are mutually consistent \eqref{eq:Xi_c}--\eqref{eq:c_Xi}.
  \end{itemize}
  \item Existence and uniqueness follow the sufficient conditions of \citet[Proposition~1]{allenarkolakis2022}. \app{appexist}
\end{itemize}
\end{frame}
\note{\textbf{Talk track, 35 seconds.} Equilibrium closes the loop among location choice, route choice, link traffic, and directed physical time. The empirical section recovers the fundamentals and anchors the congestion equation to observed physical minutes before solving counterfactual equilibria.}

% =====================================================================
\section{Data and Estimation}\label{sec:quant}
% =====================================================================
\begin{frame}[label=data]{Model Data: From Observations to Spatial Objects}
\begin{columns}[T,onlytextwidth]
\column{0.54\textwidth}
\begin{itemize}\footnotesize\setlength\itemsep{0.45em}
  \item \textbf{Residence and workplace.} Baidu 2022 home and work points are assigned to $3\times3$ km cells and aggregated to $(L_i^R,L_j^F)$.
  \item \textbf{OD flows.} The full commuting matrix supports descriptive facts and validation; the baseline inversion targets the residence and workplace marginals.
  \item \textbf{Directed network.} Road geometry, functional class, lanes, and directional speed records are mapped to adjacent grid links.
  \item \textbf{Analysis sample.} The largest connected component is retained after a detour filter removes implausible links.
\end{itemize}
\column{0.44\textwidth}
\centering\scriptsize Directional speed records\par
\includegraphics[width=\linewidth,height=0.66\textheight,keepaspectratio]{speed_eg.png}
\end{columns}
\end{frame}
\note{\textbf{Talk track, 45 seconds.} Separate observed objects from model-implied objects. Residence and workplace marginals discipline the inversion. OD data support the descriptive section and validation, while route and segment loads are generated inside the model.}

\begin{frame}{Physical Link Times Are Constructed from Segment Speeds}
\begin{columns}[T,onlytextwidth]
\column{0.49\textwidth}
\begin{itemize}\footnotesize\setlength\itemsep{0.35em}
  \item \textbf{Observed AM time.} Weekday speeds from 07:00--09:00 deliver $v_a^{obs}$.
  \item \textbf{Free-flow time.} Night speeds from 22:00--06:00 deliver $v_a^{ff}$.
  \item \textbf{Pair-specific routing.} For each adjacent grid pair $(k,l)$, a road-network path $p_{kl}$ is selected and segment minutes are aggregated:
  \vspace{-0.35em}
  \begin{equation}
    m_{kl}^{q}=60\sum_{a\in p_{kl}}\frac{\ell_a}{v_a^{q}},
    \qquad q\in\{obs,ff\}.
  \end{equation}
  \vspace{-0.65em}
  \item Here $\ell_a$ is in km and $v_a^q$ is in km/hour; 60 converts hours to minutes.
  \item Direction is retained, so $m_{kl}^{q}$ need not equal $m_{lk}^{q}$.
\end{itemize}
\column{0.49\textwidth}
\centering\scriptsize AM traversal time on adjacent grid links\par
\includegraphics[width=\linewidth,height=0.67\textheight,keepaspectratio]{app_edgeproj_adjacent_tau.png}
\end{columns}
\end{frame}
\note{\textbf{Talk track, 50 seconds.} Segment speeds are converted into physical minutes. Pair-specific routing maps the detailed road network to directed adjacent-grid links. The model then applies $t_{kl}=m_{kl}^{1/\theta}$ before constructing route and bilateral costs.}

\begin{frame}[label=estparams]{Parameters and Estimation Targets}
\small
\begin{columns}[T,onlytextwidth]
\column{0.58\textwidth}
\scriptsize
\begin{tabularx}{\linewidth}{@{}p{0.18\linewidth}p{0.16\linewidth}X@{}}
\toprule
Parameter & Value & Economic role \\
\midrule
$\theta$ & $6.83$ & response to commuting costs \\
$\alpha$ & $+0.05$ & workplace agglomeration \\
$\beta$ & $-0.10$ & residential crowding \\
$\lambda$ & $1.1$ & speed-flow elasticity \\
\bottomrule
\end{tabularx}
\column{0.40\textwidth}
\begin{itemize}\footnotesize
  \item \textbf{Matched baseline objects:} $L_i^R$ and $L_j^F$.
  \item \textbf{Measured network objects:} $m_{kl}^{obs}$, $m_{kl}^{ff}$, $n_{kl}$, and $\ell_{kl}$.
  \item \textbf{Model-implied objects:} $L_{ij}$, route shares, and $\Xi_{kl}$.
\end{itemize}
\end{columns}
\vspace{0.45em}
{\footnotesize $\theta$, $\alpha$, $\beta$, and $\lambda$ are externally calibrated in the current specification. \app{appparams}}
\end{frame}

\begin{frame}{Estimation: Invert, Anchor, and Solve}
\begin{itemize}\small
  \item \textbf{Invert fundamentals.} Recover $(\bar u_i,\bar A_j)$ so that observed residence and workplace marginals form an equilibrium at $m_{kl}^{obs}$.
  \item \textbf{Construct the baseline link-load anchor.} $\Xi_{kl}^{0}$ is the model-implied directed link load evaluated at baseline allocation and observed directed times.
  \item \textbf{Update physical link time around that operating point:}
  \vspace{-0.45em}\begin{equation}\label{eq:anchor}
    m_{kl}=\max\!\left\{
    m_{kl}^{ff},\;
    m_{kl}^{obs}
    \left[
    \frac{\Xi_{kl}}{\Xi_{kl}^{0}}
    \frac{n_{kl}^{0}}{n_{kl}}
    \right]^{\lambda}
    \right\}.
  \end{equation}
  \vspace{-0.65em}
  \item \textbf{Map and solve in levels.} Set $t_{kl}=m_{kl}^{1/\theta}$, then recompute $(m,t,\tau,\Xi,L^R,L^F,\Pi,P,\bar W)$ with fundamentals fixed.
\end{itemize}
\vspace{-0.25em}
\begin{center}\footnotesize
$\Xi_{kl}^{0}$ is a baseline model-implied load, not an observed vehicle count.
\end{center}
\end{frame}
\note{\textbf{Talk track, 55 seconds.} The inversion recovers fundamentals that rationalize the observed marginals. The anchor fixes the baseline operating point at observed physical minutes. Only the ratio $\Xi_{kl}/\Xi_{kl}^{0}$ enters, so the scale of model-implied link load cancels. Each update maps $m_{kl}$ into $t_{kl}$ before recomputing route and bilateral costs.}

\begin{frame}[label=fundamentals]{Baseline Fundamentals (inverted)}
\begin{columns}[T,onlytextwidth]
\column{0.5\textwidth}
\centering \scriptsize Residence amenity $\bar u_i$\\
\includegraphics[width=\linewidth,height=0.52\textheight,keepaspectratio]{fund_ubar.png}
\column{0.5\textwidth}
\centering \scriptsize Workplace productivity $\bar A_j$\\
\includegraphics[width=\linewidth,height=0.52\textheight,keepaspectratio]{fund_abar.png}
\end{columns}
\vspace{0.1em}
\begin{itemize}\footnotesize
  \item $\bar u,\bar A$ are inverted to match observed residence and workplace shares at observed travel times.
  \item $\mathrm{corr}(\bar A_j,L_j^F)=\good{0.67}$ vs.\ $\mathrm{corr}(\bar u_i,L_i^R)=0.30$: jobs pinned by productivity, residence is access-driven.
  \item Treated corridors land on \bad{high-flow employment bottlenecks} (TidalCandidateScore vs.\ slow-direction $\Xi_s$: corr.\ $\approx+0.91$ on selected top-100 pairs). \app{appfund}
\end{itemize}
\end{frame}
\note{\textbf{Talk track, 35 seconds.} These maps are recovered objects, not directly observed amenities or productivity. The inversion matches baseline residence and workplace marginals at observed directed times. Use the maps to establish the spatial fundamentals that remain fixed in the tidal-lane counterfactuals.}

% =====================================================================
\section{Counterfactuals}
% =====================================================================
\begin{frame}{What Symmetry Misses}
\framesubtitle{The directional and symmetric measurement systems are re-estimated separately.}
\small
\begin{center}
\scriptsize
\begin{tabularx}{0.94\textwidth}{@{}p{0.18\textwidth}p{0.25\textwidth}p{0.24\textwidth}X@{}}
\toprule
World & Link-time measurement & Fundamental inversion & Equilibrium \\
\midrule
Directional & Directed $m_{kl}^{obs}$ and $m_{kl}^{ff}$ & Re-estimate $(\bar u,\bar A)$ & Re-solve anchored GE \\
Symmetric & Pairwise geometric mean & Re-estimate $(\bar u,\bar A)$ & Re-solve anchored GE \\
\bottomrule
\end{tabularx}
\end{center}
\vspace{0.25em}
\begin{itemize}\footnotesize
  \item \textbf{Primary estimand.} Changes in recovered fundamentals, market access, and equilibrium spatial gradients induced by imposing reciprocal symmetry.
  \item \textbf{Distributional diagnostics.} Compare dispersion, CBD gradients, and dissimilarity in $\bar u$, $\bar A$, bilateral costs, and model-implied link loads.
  \item \textbf{Policy relevance.} Re-evaluate the tidal-lane counterfactual under each re-estimated measurement system.
\end{itemize}
\vspace{0.3em}
\begin{center}
\fcolorbox{gray!50}{gray!8}{\parbox{0.77\textwidth}{\centering\footnotesize
[Full re-estimation results pending.] The fixed-fundamentals diagnostic is retained in the appendix.}}
\end{center}
\end{frame}

\begin{frame}{Where the Tidal Policy Lands}
\framesubtitle{Top-100 corridors by TidalCandidateScore (TCS).}
\begin{columns}[T,totalwidth=0.96\textwidth]
\column{0.42\textwidth}
\begin{itemize}\footnotesize
  \item For each reciprocal pair, label the slower AM direction $s$ and the faster direction $f$.
  \item Reallocate one lane while holding total road width fixed:
  \vspace{-0.35em}
  \[
    n_s'=n_s+1,\qquad n_f'=n_f-1.
  \]
  \vspace{-0.75em}
  \item \textbf{TCS} is the slow-minus-fast congestion-delay mass:
  \vspace{-0.35em}\[
    \begin{aligned}
    \mathrm{TCS}={}&\Xi_s(m_s^{obs}-m_s^{ff})\\
    &-\Xi_f(m_f^{obs}-m_f^{ff}).
    \end{aligned}
  \]
  \vspace{-0.65em}
  \item Red edges receive the lane among the top-100 reciprocal pairs.
\end{itemize}
\column{0.54\textwidth}
\centering
\vspace{-0.2em}
\includegraphics[width=\linewidth,height=0.78\textheight,keepaspectratio]{map_treated_overlay.png}
\end{columns}
\end{frame}
\note{\textbf{Talk track, 60 seconds.} The intervention is capacity reallocation, not capacity expansion. One lane moves from the faster AM direction to the slower direction on each selected reciprocal corridor. TCS ranks corridors by the difference in baseline flow times excess traversal time. The map identifies the lane-receiving directions.}

\begin{frame}{Short Run vs.\ Long Run}
\small
\begin{itemize}
  \item \textbf{Short run (ST):} pin residence $L^R$ and employment $L^F$ at baseline; routes, link traffic, and physical link times $m_{kl}$ re-equilibrate.
  \item \textbf{Long run (LT):} allow $L^R$, $L^F$, and OD shares to sort, with congestion feedback.
\end{itemize}
\vspace{-0.25em}
{\scriptsize
\begin{equation*}
\bar M^h
=\frac{1}{\bar L}\sum_{k,l}\Xi_{kl}^h m_{kl}^h
=\frac{1}{\bar L}\sum_{i,j,r}L_{ij,r}^h
  \sum_{(k,l)\in r}m_{kl}^h,
\qquad h\in\{0,ST,LT\}.
\end{equation*}
}
\vspace{-0.55em}
\begin{itemize}\footnotesize\setlength\itemsep{0.35em}
  \item $\bar M^h$ is expected additive physical time; seconds saved per commuter equal $-60(\bar M^h-\bar M^0)$.
  \item The bilateral index $\tau_{ij}$ remains in choice and welfare equations, but is never converted into minutes or seconds. \app{appphysicaltime}
\end{itemize}
\end{frame}
\note{\textbf{Talk track, 60 seconds.} The short run fixes residence and workplace marginals while route traffic and congestion adjust. The long run also permits sorting. Physical time is computed from equilibrium link passages times physical link minutes, which equals the route-level sum of minutes under the model. The composite bilateral index $\tau_{ij}$ remains a choice and welfare object, not a clock-time measure.}

\begin{frame}[label=result2]{Short-Run Time Gains Are Reversed by Long-Run Reallocation}
\framesubtitle{Citywide expected physical time from equilibrium link passages and minutes.}
\centering
\begin{minipage}[t]{0.62\textwidth}
\centering
\small\textbf{Physical commuting time}\par
\vspace{0.35em}
\scriptsize
\renewcommand{\arraystretch}{1.18}
\setlength{\tabcolsep}{6pt}
\begin{tabular}{lcc}
\toprule
 & Short run & Long run \\
\midrule
Average minutes & $10.60$ & $11.54$ \\
Change, seconds/commuter & \good{$-42.1$} & \bad{$+14.4$} \\
Person-hours saved, AM peak & \good{$+1{,}031$} & \bad{$-352$} \\
\bottomrule
\end{tabular}
\end{minipage}\hfill
\begin{minipage}[t]{0.30\textwidth}
\centering
\small\textbf{Expected-utility welfare}\par
\vspace{0.35em}
\scriptsize
\begin{tabular}{lc}
\toprule
Equilibrium & $\Delta\bar W$ \\
\midrule
Short run & \good{$+0.87\%$} \\
Long run & \good{$+0.32\%$} \\
\bottomrule
\end{tabular}
\end{minipage}

\vspace{0.60em}
\begin{itemize}\footnotesize\setlength\itemsep{0.35em}
  \item \textbf{Baseline.} Average physical commuting time is $11.30$ minutes.
  \item \textbf{Adjustment.} The $42.1$-second ST gain is offset by $26.6$ seconds of additional network feedback and $30.0$ seconds from long-run link-passage reallocation.
  \item \textbf{Welfare.} Expected utility remains above baseline because long-run choices also reflect productivity, amenities, wages, and access. Mode switching remains outside the model.
\end{itemize}
\end{frame}
\note{\textbf{Talk track, 90 seconds.} The policy lowers citywide expected physical commuting time from 11.30 to 10.60 minutes in the short run, a gain of 42.1 seconds per commuter or about 1,031 person-hours in one AM peak. After long-run sorting and congestion feedback, average time reaches 11.54 minutes, 14.4 seconds above baseline. Expected utility remains positive because households jointly value wages, productivity, amenities, and access. Physical time and welfare therefore need not move in the same direction.}

\begin{frame}{Long-Run Sorting}
\begin{columns}[T,onlytextwidth]
\column{0.35\textwidth}
\begin{itemize}\footnotesize
  \item Totals are fixed at $\bar L$; maps report demeaned share changes, not citywide growth.
  \item Positive cells gain share relative to baseline; black circles mark treated-corridor endpoint nodes.
  \item \textbf{Residential shares} rise near treated endpoints because lower OD costs raise residential market access.
  \item \textbf{Employment gains} are concentrated in a smaller set of red cells.
\end{itemize}
\column{0.63\textwidth}
\centering
{\scriptsize
\begin{tabular}{lcccc}
\toprule
 & \multicolumn{2}{c}{Treated endpoints} & \multicolumn{2}{c}{Rest} \\
 & $\ell^R$ & $\ell^F$ & $\ell^R$ & $\ell^F$ \\
\midrule
Mean $\Delta\log$ & \good{$+0.017$} & $-0.006$ & \bad{$-0.012$} & \bad{$-0.016$} \\
\bottomrule
\end{tabular}
}

\vspace{1em}
\begin{minipage}[t]{0.49\linewidth}
\centering \scriptsize demeaned $\Delta\log\ell^R$\\
\includegraphics[width=\linewidth,height=0.45\textheight,keepaspectratio]{sort_lR.png}
\end{minipage}\hfill
\begin{minipage}[t]{0.49\linewidth}
\centering \scriptsize demeaned $\Delta\log\ell^F$\\
\includegraphics[width=\linewidth,height=0.45\textheight,keepaspectratio]{sort_lF.png}
\end{minipage}
\end{columns}
\end{frame}

\begin{frame}[label=result4]{Transportation Policy Comparison}
\framesubtitle{Endogenous GE under anchored congestion and long-run sorting.}
\begin{center}\scriptsize
\begin{tabular}{lcccc}
\toprule
Policy & ST $\bar W$ & LT $\bar W$ & net lane-km & cost/lane-km \\
\midrule
Tidal reallocation & $+0.87\%$ & $+0.32\%$ & \good{\textbf{0}} & near zero capex \\
Road Expansion & $+1.26\%$ & $+0.44\%$ & 1{,}140 & $\approx$ CNY 5m \\
\bottomrule
\end{tabular}
\end{center}
\vspace{0.3em}
\begin{itemize}\footnotesize
  \item Tidal reallocation lifts LT welfare $+0.32\%$ at zero net lane-km.
  \item Delivering comparable long-run welfare by building capacity takes about $1{,}140$ lane-km: Road Expansion adds lanes where baseline flow$\times$delay is highest and reaches $+0.44\%$.
  \item Urban arterial widening runs on the order of CNY 5m per lane-km (roughly CNY 30m/km for a six-lane arterial), so the expansion benchmark is order CNY 6bn; tidal reallocation avoids this outlay. \app{apppolicy}
\end{itemize}
\end{frame}
\note{\textbf{Talk track, 45 seconds.} The comparison is limited to welfare, capacity, and construction cost because the two policies do not share the same treated-link support. Tidal reallocation changes the direction in which existing capacity is deployed; the expansion benchmark builds about 1,140 lane-kilometers. The preceding slide reports citywide physical time separately from expected utility.}

\begin{frame}{Conclusion}
\begin{enumerate}
  \item \textbf{Directional congestion is a widespread urban feature.}
  \begin{itemize}\footnotesize
    \item Across the top 100 congested cities, flow-weighted median \SAR\ is \good{0.878}; the same asymmetry persists on Beijing road segments.
  \end{itemize}
  \item \textbf{Directional measurement belongs inside the structural inversion.}
  \begin{itemize}\footnotesize
    \item The full comparison re-estimates fundamentals and equilibrium gradients under the directional and reciprocal-symmetric cost surfaces.
  \end{itemize}
  \item \textbf{Tidal lanes deliver welfare gains through directed capacity reallocation.}
  \begin{itemize}\footnotesize
    \item Physical time falls by \good{$42.1$ seconds per commuter} in ST but rises by \bad{$14.4$ seconds} above baseline after LT reallocation; $\bar W$ remains \good{$+0.32\%$} in LT.
    \item Supply-side policies with comparable LT welfare require about \bad{$1{,}100$ lane-km}; tidal reallocation uses \good{zero net lane-km}.
  \end{itemize}
\end{enumerate}
\end{frame}
\note{\textbf{Talk track, 35 seconds.} Directional congestion is widespread across cities and persists at the road-segment level. A directed spatial model connects this asymmetry to route traffic, congestion feedback, and sorting. Tidal reallocation produces a clear short-run physical-time gain, while long-run reallocation reverses the average-time effect. Expected utility remains above baseline, and the policy uses zero net lane-kilometers.}

% =====================================================================
\appendix
\section{Appendix}
% =====================================================================
\begin{frame}[label=appsymfixed]{Appendix: Fixed-Fundamentals Symmetry Diagnostic}
\framesubtitle{This estimand differs from the full re-estimation in the main deck.}
\footnotesize
Hold the network, lanes, inverted fundamentals, parameters, solver, and welfare formula fixed; replace reciprocal-pair $m_{kl}^{obs}$ and $m_{kl}^{ff}$ by pairwise geometric means before applying $t_{kl}=m_{kl}^{1/\theta}$.
\vspace{0.45em}
\begin{center}
\scriptsize
\begin{tabular}{lcccc}
\toprule
 & $\Delta\bar W$ & Residence & Workplace & OD cost \\
 &  & dissim. & dissim. & $\Delta$ \\
\midrule
Symmetric vs.\ directed & \good{$+1.10\%$} & $3.36\%$ & $6.15\%$ & $+0.10\%$ \\
\bottomrule
\end{tabular}
\end{center}
\vspace{0.35em}
\begin{itemize}\footnotesize
  \item The experiment changes the measured cost surface while holding fundamentals fixed.
  \item It is not a comparison of two fully re-estimated models.
\end{itemize}
\end{frame}

\begin{frame}[label=approute]{Appendix: Route Costs and Bilateral Commuting Costs}
\begin{itemize}\small
  \item Define the directed adjacency matrix $A=[a_{kl}]$ by
  \vspace{-0.5em}\begin{equation}
    a_{kl} = t_{kl}^{-\theta}
  \end{equation}
  \item Summing over all feasible route lengths gives
  \vspace{-0.5em}\begin{equation}
    B = \sum_{K=0}^{\infty}A^K = (I-A)^{-1}
  \end{equation}
  \item The bilateral commuting-cost aggregator can be written as
  \vspace{-0.5em}\begin{equation}\label{eq:inverse}
    \tau_{ij} = B_{ij}^{-1/\theta}
  \end{equation}
  \vspace{-0.5em}
  \item The implementation applies a common link-cost rescaling $\zeta$ to keep the spectral radius of $A$ below one. This scalar is distinct from $\kappa_{ij}^{kl}$, the expected use of link $(k,l)$ by an $i\to j$ commuter.
\end{itemize}
\navback{bilateral}
\end{frame}

\begin{frame}[label=appphysicaltime]{Appendix: Expected Physical Travel Time}
\scriptsize
\begin{itemize}\setlength\itemsep{0.45em}
  \item Physical minutes add along route $r$:
  \vspace{-0.35em}
  \[
    M_{ij,r}=\sum_{(k,l)\in r}m_{kl}.
  \]
  \vspace{-0.75em}
  \item Iceberg link costs compound, $C_{ij,r}=\prod_{(k,l)\in r}t_{kl}$, and the conditional route probability is
  \vspace{-0.35em}
  \[
    \rho_{ij,r}=\frac{C_{ij,r}^{-\theta}}{\sum_{s\in\mathcal R_{ij}}C_{ij,s}^{-\theta}},
    \qquad t_{kl}=m_{kl}^{1/\theta}.
  \]
  \vspace{-0.75em}
  \item Expected OD minutes and the citywide mean are
  \vspace{-0.35em}
  \[
    M_{ij}=\sum_r\rho_{ij,r}M_{ij,r},
    \qquad
    \bar M=\frac{1}{\bar L}\sum_{i,j}L_{ij}M_{ij}
    =\frac{1}{\bar L}\sum_{k,l}\Xi_{kl}m_{kl}.
  \]
  \vspace{-0.75em}
  \item No transformation of $\tau_{ij}$ produces additive route minutes. The implementation computes $\bar M$ from equilibrium link passages and $m_{kl}$.
\end{itemize}
\navback{result2}
\end{frame}

\begin{frame}[label=apptreatedtime]{Appendix: Treated-Link Time Incidence}
\framesubtitle{A local engineering incidence measure, distinct from citywide experienced travel time.}
\begin{columns}[T,onlytextwidth]
\column{0.52\textwidth}
\scriptsize
\begin{align*}
\text{Gross savings}^{h}
  &=\sum_{e\in\mathcal E^*}\max\{m_e^0-m_e^h,0\}\,\Xi_e^0,\\
\text{Offsetting losses}^{h}
  &=\sum_{e\in\mathcal E^*}\max\{m_e^h-m_e^0,0\}\,\Xi_e^0.
\end{align*}
\begin{itemize}\scriptsize
  \item $\mathcal E^*$ contains the 200 treated directed links.
  \item Baseline link-passage weights isolate time changes on treated links.
\end{itemize}
\column{0.45\textwidth}
\centering\scriptsize
\begin{tabular}{lcc}
\toprule
Person-hours & ST & LT \\
\midrule
Gross savings & $1{,}385$ & $804$ \\
Offsetting losses & $-109$ & $-90$ \\
Net savings & $1{,}276$ & $714$ \\
\bottomrule
\end{tabular}
\vspace{0.55em}

\raggedright\scriptsize A commuter can traverse multiple treated links. These totals therefore cannot be divided by unique exposed commuters to obtain seconds per commuter.
\end{columns}
\navback{result2}
\end{frame}

\begin{frame}[label=appfund]{Appendix: Estimation Algorithm}
\footnotesize
\textbf{Baseline inversion}
\begin{itemize}
  \item Treat observed residence and workplace marginals as the baseline equilibrium.
  \item Recover $(\bar u_i^\theta,\bar A_j^\theta)$ by alternating updates with geometric-mean normalization.
  \item Match $(L_i^R,L_j^F)$, not the full bilateral OD matrix; OD flows are model-implied by gravity.
\end{itemize}
\vspace{0.25em}
\textbf{Counterfactual solve}
\begin{itemize}
  \item Given fundamentals and parameters, solve $(L^R,L^F)$, market access, model-implied link load $\Xi$, and physical link times jointly.
  \item Anchored congestion updates physical time around the baseline model-implied load and observed lanes:
  \vspace{-0.3em}
  \[
    m_{kl}=\max\{m_{kl}^{ff},\,m_{kl}^{obs}[(\Xi_{kl}/\Xi_{kl}^{0})(n_{kl}^{0}/n_{kl})]^\lambda\},
    \qquad t_{kl}=m_{kl}^{1/\theta}.
  \]
  \vspace{-0.8em}
  \item Outer loop updates $m_{kl}$; inner loop alternates location shares and market access until convergence.
\end{itemize}
\navback{fundamentals}
\end{frame}

\begin{frame}[label=appduncan]{Appendix: Jobs--Housing Dissimilarity Benchmarks}
\framesubtitle{Duncan $D=\tfrac12\sum_g\lvert j_g/J-r_g/R\rvert$ ($j,r=$ jobs, residents per unit $g$); aspatial, $0$ balanced $\to1$ separated.}
\begin{columns}[T,onlytextwidth]
\column{0.60\textwidth}
\centering
\includegraphics[width=\linewidth,height=0.58\textheight,keepaspectratio]{app_duncan_cross_city.png}
\column{0.38\textwidth}
\begin{itemize}\footnotesize
  \item \textbf{Cross-city.} Township $D$, top-100 cities (residents $=$ origin, jobs $=$ dest mass, 2021 flows): Beijing $\good{0.20}$ ranks \good{1/100} (median $0.09$).
  \item \textbf{Caveat.} $D$ is aspatial, rising with size (corr $0.50$; $0.21$ with unit count): Beijing's lead is not a unit-count artefact, but $D$ ignores distance.
\end{itemize}
\end{columns}
\navback{jobhousing}
\end{frame}

\begin{frame}[label=appinbound]{Appendix: Radial AM Inbound--Reverse Asymmetry}
\framesubtitle{Road-mode OD; \emph{inbound} $=$ home $\geq15$ km from the job-weighted center \& work $\leq10$ km; \emph{reverse} $=$ mirror.}
\begin{columns}[T,onlytextwidth]
\column{0.62\textwidth}
\centering
\includegraphics[width=\linewidth,height=0.66\textheight,keepaspectratio]{cityShenzhen2.jpg}
\column{0.36\textwidth}
\begin{itemize}\footnotesize
  \item Each core-directed trip is binned into one of 16 $22.5^\circ$ sectors by \emph{home} bearing.
  \item Pooled inbound share \good{$74.5\%$}; inbound/reverse $\good{2.9\times}$; \good{11/16} sectors $>0.70$ (min $0.57$, median $0.77$, max $0.86$).
  \item Non-circular: inbound and reverse share one center, so the ratio is a genuine imbalance, not an artefact of the job-cluster construction.
\end{itemize}
\end{columns}
\navback{jobhousing}
\end{frame}

\begin{frame}[label=appmodelasym]{Appendix: Model-Implied Directional Asymmetry}
\framesubtitle{Baseline directed traffic $\Xi$: ratio of the two directions per link (model output, not raw data).}
\setlength{\columnsep}{0.005\textwidth}
\begin{columns}[T,onlytextwidth]
\column{0.53\textwidth}
\centering
\includegraphics[width=\linewidth,height=0.70\textheight,keepaspectratio]{map_dir_asym.png}
\column{0.45\textwidth}
\begin{itemize}\footnotesize
  \item Color is $\log_{10}(\Xi_{\mathrm{high}}/\Xi_{\mathrm{low}})$; width is total reciprocal flow.
  \item Median two-way imbalance is \bad{$2.9\times$}; p90 is \bad{$20\times$}.
  \item High-imbalance links concentrate on the central high-flow spine.
  \item Object: equilibrium link load $\Xi$, not observed speed.
\end{itemize}
\end{columns}
\navback{equilibrium}
\end{frame}

\begin{frame}[label=appparams]{Appendix: Parameter Interpretation}
\begin{itemize}\small
  \item $\theta$ governs sensitivity to bilateral commuting costs. Higher $\theta$ makes households and jobs respond more sharply to directed cost changes.
  \item $\alpha>0$ is workplace agglomeration in $w_j=\bar A_j(L_j^F)^\alpha$. It reinforces employment concentration around locations that gain access.
  \item $\beta<0$ is residential crowding in $u_i=\bar u_i(L_i^R)^\beta$. It penalizes residential pile-up and dampens household concentration.
  \item $\lambda$ maps traffic per lane into directed physical link time. Higher $\lambda$ raises the short-run congestion response to capacity reallocation.
  \item Sorting reflects the interaction of all four forces: accessibility changes shift OD costs, $\alpha$ pulls jobs toward stronger access, and $\beta$ limits residential concentration.
\end{itemize}
\navback{estparams}
\end{frame}

\begin{frame}[label=appsens]{Appendix: Sensitivity Summary}
\framesubtitle{Final mainline rerun, TCS top-100, primitive $\bar W$, $\lambda=1.1$.}
\begin{columns}[T,onlytextwidth]
\column{0.5\textwidth}
\centering Vary $\alpha$ ($\beta=-0.10$)\\[2pt]
\scriptsize
\setlength{\tabcolsep}{10pt}
\begin{tabular}{cc}
\toprule
$\alpha$ & $\bar W$ LT \\
\midrule
$-0.12$ & $+0.37\%$ \\
$0$     & $+0.34\%$ \\
$+0.03$ & $+0.33\%$ \\
$\mathbf{+0.05}$ & $\mathbf{+0.32\%}$ \\
$+0.10$ & $+0.29\%$ \\
\bottomrule
\end{tabular}
\column{0.5\textwidth}
\centering Vary $\beta$ ($\alpha=+0.05$)\\[2pt]
\scriptsize
\setlength{\tabcolsep}{10pt}
\begin{tabular}{cc}
\toprule
$\beta$ & $\bar W$ LT \\
\midrule
$\mathbf{-0.10}$ & $\mathbf{+0.32\%}$ \\
$0$     & $+0.31\%$ \\
$+0.05$ & $+0.28\%$ \\
$+0.10$ & \emph{not converged} \\
\bottomrule
\end{tabular}
\end{columns}
\vspace{0.25em}
\begin{itemize}\footnotesize
  \item $\bar W$ LT remains positive in every converged $\alpha$--$\beta$ cell; route-time statistics move more because sorting changes OD weights.
  \item Detailed cells, $\theta$ grid, and congestion-band checks\app{appsensdetail}\app{appsens_theta}.
\end{itemize}
\navback{result2}
\end{frame}

\begin{frame}[label=appsensdetail]{Appendix: Sensitivity Details}
\scriptsize
\begin{columns}[T,onlytextwidth]
\column{0.40\textwidth}
\begin{flushright}
\setlength{\tabcolsep}{4pt}
\begin{tabular}{ccc}
\toprule
$\alpha$ & $\beta$ & $\bar W$ LT \\
\midrule
$-0.12$ & $-0.10$ & $+0.37\%$ \\
$0.00$ & $-0.10$ & $+0.34\%$ \\
$0.03$ & $-0.10$ & $+0.33\%$ \\
$\mathbf{0.05}$ & $\mathbf{-0.10}$ & $\mathbf{+0.32\%}$ \\
$0.10$ & $-0.10$ & $+0.29\%$ \\
$0.05$ & $0.00$ & $+0.31\%$ \\
$0.05$ & $0.05$ & $+0.28\%$ \\
$0.05$ & $0.10$ & \emph{not converged} \\
\bottomrule
\end{tabular}
\end{flushright}
\column{0.58\textwidth}
\begin{itemize}\footnotesize
  \item Mainline ST $\bar W$ is $+0.87\%$; citywide physical time is reported on Slide~\ref{result2}, while treated-link incidence is retained in Appendix~\ref{apptreatedtime}.
  \item $\theta$ grid (\app{appsens_theta}): each value re-inverts fundamentals and GE; LT $\bar W$ gains decline with $\theta$ while signs stay positive.
  \item Congestion band: at $\lambda=0.57$, top-100 LT $\bar W$ is $-0.078\%$; this is retained as a legacy lower-slope robustness, not the mainline.
  \item Selection robustness: the asymmetry-ratio rule is weaker than TCS selection and can flip LT $\bar W$ negative.
\end{itemize}
\end{columns}
\navback{appsens}
\end{frame}

\begin{frame}[label=appsens_theta]{Appendix: Theta Sensitivity}
\framesubtitle{TCS top-100; each $\theta$ re-inverts fundamentals and re-solves baseline + tidal ST/LT.}
\scriptsize
\begin{center}
\renewcommand{\arraystretch}{1.3}
\setlength{\tabcolsep}{12pt}
\begin{tabular}{ccccc}
\toprule
$\theta$ & ST $\bar W$ & LT $\bar W$ & res.\ dissim & job dissim \\
\midrule
$3$ & $+1.99\%$ & $+0.75\%$ & $1.11$ & $1.55$ \\
$5$ & $+1.19\%$ & $+0.45\%$ & $0.97$ & $1.67$ \\
$\mathbf{6.83}$ & $\mathbf{+0.87\%}$ & $\mathbf{+0.32\%}$ & $\mathbf{0.86}$ & $\mathbf{1.79}$ \\
$8$ & $+0.74\%$ & $+0.27\%$ & $0.80$ & $1.89$ \\
$10$ & $+0.59\%$ & $+0.21\%$ & $0.70$ & $2.07$ \\
\bottomrule
\end{tabular}
\end{center}
\vspace{0.4em}
\begin{itemize}\footnotesize
  \item LT $\bar W$ stays positive across the grid; magnitudes fall as $\theta$ rises (stronger cost sensitivity).
  \item Resident dissim falls and job dissim rises with $\theta$; baseline $\bar W$ levels are not comparable across $\theta$.
\end{itemize}
\note{Mainline $\theta=6.83$ matches the tidal diagnostics row; baseline $\bar W$ levels are not comparable across $\theta$.}
\navback{appsens}
\end{frame}

\begin{frame}[label=apppolicy]{Appendix: Policy Definitions and Selected Effects}
\scriptsize
\begin{center}
\resizebox{0.98\textwidth}{!}{%
\begin{tabular}{p{0.19\textwidth}p{0.39\textwidth}cccc}
\toprule
Policy & Capacity rule & ST $\bar W$ & LT $\bar W$ & lane-km & cost/lane-km \\
\midrule
Tidal reallocation & TCS top-100; add one lane in the slow direction and remove one from the reciprocal fast direction. & $+0.87\%$ & $+0.32\%$ & 0 & near zero capex \\
Road Expansion & Top 100 by baseline flow $\times$ excess travel time; add one lane to each direction. & $+1.26\%$ & $+0.44\%$ & 1{,}140 & $\approx$ CNY 5m \\
Uniform expansion & Scale all links by $+2.4\%$ so added lane-km matches the targeted-expansion budget. & $+0.21\%$ & $+0.10\%$ & 1{,}092 & $\approx$ CNY 5m \\
\bottomrule
\end{tabular}
}
\end{center}
\vspace{0.2em}
\begin{itemize}\footnotesize
  \item ST pins residence and workplace shares at baseline; LT re-solves residence--workplace sorting.
  \item Road Expansion asks where an engineer would add capacity using flow and delay only; uniform spreads the same lane-km budget across the whole network (targeting versus spreading).
  \item Cost benchmark uses 2023--24 municipal urban-arterial construction estimates ($\approx$ CNY 30m per km, roughly CNY 5m per widened lane-km, a single widened lane being a fraction of the full cross-section), not inter-city expressway costs; the targeted expansion is thus order CNY 6bn.
\end{itemize}
\navback{result4}
\end{frame}

\begin{frame}[label=resultelast]{Appendix: Targeting Tidal Lanes}
\framesubtitle{Per-link welfare elasticity from $d\bar W^\theta=dA_{ab}\phi_a\psi_b$.}
\begin{columns}[T,onlytextwidth]
\column{0.44\textwidth}
\centering
\includegraphics[width=\linewidth,height=0.66\textheight,keepaspectratio]{map_link_welfare_elast.png}
\column{0.53\textwidth}
\scriptsize
\textbf{Selection rule} (top-100):
\begin{tabular}{lc}
\toprule
rule & LT $\bar W$ \\
\midrule
\good{welfare elast.} & \good{$+0.35\%$} \\
TCS score & $+0.32\%$ \\
speed only & $+0.13\%$ \\
\bottomrule
\end{tabular}
\vspace{0.5em}
\begin{itemize}\footnotesize
  \item The local score decomposes into edge gain $\times$ upstream residence access $\times$ downstream workplace opportunity. \app{appelast}
  \item Flow matters: the speed-only rank correlation with welfare elasticity is \bad{$0.08$}; the flow$\times$delay correlation is $0.65$.
  \item Welfare-elasticity and TCS targeting select the high-flow spine; speed-only also selects slow, low-flow links.
\end{itemize}
\end{columns}
\end{frame}

\begin{frame}[label=appelast]{Appendix: Welfare-Elasticity Score}
\scriptsize
\begin{itemize}
  \item Define the directed-link accessibility matrix by $A_{ab}=t_{ab}^{-\theta}$ and the network multiplier by $B=(I-A)^{-1}$.
  \item A local gain on one directed edge uses the matrix-inverse differential:
  \[
    dB=B(dA)B.
  \]
  \item If only edge $a\to b$ changes, $dA$ has one nonzero entry, so the welfare-index numerator has the local derivative
  \[
    d\bar W^\theta=dA_{ab}\phi_a\psi_b.
  \]
  \item $dA_{ab}$ is the direct gain from one more lane; $\phi_a$ is upstream residence-access weight reaching node $a$; $\psi_b$ is downstream workplace-opportunity weight from node $b$.
  \item The score is a local ranking device. The reported LT $\bar W$ values still come from re-solving the endogenous GE counterfactual.
\end{itemize}
\navback{resultelast}
\end{frame}

\begin{frame}[label=nextsteps]{Appendix: Next Steps}
\small
\begin{itemize}
  \item \textbf{Reduced-form evidence on tidal lanes.}
  \begin{itemize}\footnotesize
    \item Record traffic at operational tidal-lane segments and recover directed vehicle counts and speeds.
    \item Estimate the local speed--flow relation around the lane-direction switch. \app{appfieldvisit}
  \end{itemize}
  \item \textbf{Structural parameters.}
  \begin{itemize}\footnotesize
    \item Tighten estimates of $\alpha$, $\beta$, $\theta$, and $\lambda$ as the data permit.
  \end{itemize}
  \item \textbf{External validity.}
  \begin{itemize}\footnotesize
    \item Extend the measurement and counterfactual pipeline to Beijing and Shenzhen, AM and PM peaks, and alternative grid resolutions.
  \end{itemize}
\end{itemize}
\end{frame}

\begin{frame}[label=appfieldvisit]{Appendix: Field Visit}
\framesubtitle{On-site survey of Shenzhen tidal-lane corridors.}
\vspace{-0.15em}
\setlength{\columnsep}{0.014\textwidth}
\begin{columns}[T,onlytextwidth]
\column{0.317\textwidth}
\centering
\includegraphics[width=\linewidth,height=0.305\textheight,keepaspectratio]{fieldvisit_crop_1.jpg}\par\vspace{0.24em}
\includegraphics[width=\linewidth,height=0.305\textheight,keepaspectratio]{fieldvisit_crop_2.jpg}
\column{0.317\textwidth}
\centering
\includegraphics[width=\linewidth,height=0.305\textheight,keepaspectratio]{fieldvisit_crop_5.jpg}\par\vspace{0.24em}
\includegraphics[width=\linewidth,height=0.305\textheight,keepaspectratio]{fieldvisit_crop_4.jpg}
\column{0.317\textwidth}
\centering
\includegraphics[width=\linewidth,height=0.305\textheight,keepaspectratio]{fieldvisit_crop_6.jpg}\par\vspace{0.24em}
\includegraphics[width=\linewidth,height=0.305\textheight,keepaspectratio]{fieldvisit_crop_7.jpg}
\end{columns}
\navback{nextsteps}
\end{frame}

\begin{frame}[label=appexist]{Appendix: Existence and Uniqueness}
\framesubtitle{Urban commuting model; sufficient conditions from \citet[Proposition~1]{allenarkolakis2022}.}
\small
\begin{itemize}
  \item \textbf{Existence.} On a strongly connected network a strictly positive equilibrium exists for any $\alpha,\beta$, $\theta>0$, and $\lambda\ge0$ (Brouwer fixed point on the transformed operator).
  \item \textbf{Uniqueness (sufficient), for $\alpha,\beta\in[-1,1]$:}
  \vspace{-0.15em}\[
    \alpha\le\tfrac12\!\left(\tfrac1\theta-\lambda\right)
    \qquad\text{and}\qquad
    \beta\le\tfrac12\!\left(\tfrac1\theta-\lambda\right).
  \]
  \item \textbf{Principle.} The proof bounds the equilibrium operator so that its elasticity map has spectral radius at most one \citep{allenarkolakisli2020,karlin1967}.
  \item \textbf{Intuition.} $\alpha$ and $\beta$ govern agglomeration forces; stronger congestion $\lambda$ tightens the sufficient uniqueness condition.
  \item The condition is sufficient, not necessary. At $(\alpha,\beta)=(0.05,\,0.10)$ both spillovers turn agglomerative, so the iteration fails to converge.
\end{itemize}
\navback{equilibrium}
\end{frame}

\nobibliography{references}

\end{document}
